%======================================================================
\section{The certificates}\label{app:cert}
%======================================================================

This appendix describes the four computer-assisted ingredients listed in Section~\ref{ssec:claims}. Each has a standalone script in the
ancillary files, with its exact inputs pinned by SHA-256, the command that was run, and the log of a fresh run; the file \ancf{README.md} of
the ancillary files lists the expected decisive output lines and the runtimes. Paths below are relative to the top directory of the
ancillary files. In the scripts and logs the function $\Pi$ of Theorem~\ref{thm:F} is called \texttt{R}, $\Ax$ is \texttt{Phi}, $\Bx$ is \texttt{g},
$\ck_t$ is \texttt{fhat\_t}; the large-$|t|$ theorem, the window and the moment certificate are called Theorem~L, Certificate~W and
Certificate~M (Certificate~M-box on the trivial-bound box $\TB$). Certificate~W and the run of Theorem~L at $t=8$ are independent checks, not
used in the proofs of the main results.

\emph{Conventions.} All bounds are computed in ball arithmetic (Arb \cite{Joh17}, through python-flint, version~0.9.0, built on FLINT
\cite{FLINT}): every quantity is a midpoint--radius interval with outward rounding that contains the exact value, and every truncation, tail and
quadrature error is bounded explicitly. Binary64 numbers are used only to choose parameters (cell centres, partitions, ellipse parameters),
which are then used as the exact binary numbers they are. The modified Bessel function $K_0$ for real arguments is taken from the validated
library of Part~I (\ancf{lib/besselk.py}, shipped byte-identical); Arb's own \texttt{bessel\_k} gives correct but very wide enclosures at
moderate arguments and is not used for certified quantities. For complex arguments only the majorant $|K_\nu(w)|\le K_\nu(\Re w)$ ($\Re w>0$,
$\nu=0,1$) is used, which follows from $K_\nu(w)=\int_0^\infty e^{-w\cosh s}\cosh(\nu s)\dd s$.

\subsection{The coefficient enclosures}\label{app:cert-coeffs}

\emph{Statement.} Certified balls for $a_n=n\Ss_n+\frac14\delta_{n1}$: for $n\le1500$ from the terms $c\le10^4$ ($n\le60$) or
$c\le\max(600,\lceil4\pi\sqrt n\rceil+1)$ ($60<n\le1500$); for $n\le6$ from $c\le3\cdot10^5$; for $n\le2$ from $c\le6\cdot10^5$; in each case plus a
rigorous bound for the remaining terms. Certificate~W and the run of Theorem~L at $t=8$ use $n\le300$ and the bound $|a_n|\le\abar_n$ beyond; with
$c\le3\cdot10^5$,
\[
a_1=4581.370553408\pm6.9\cdot10^{-5},\qquad a_2=1025727.234918\pm2.7\cdot10^{-4},
\]
$a_3=62590223.689992\pm5.9\cdot10^{-4}$,
and every $a_n$, $n\le1500$, is certified positive, in agreement with Theorem~\ref{thm:an-positive}.

\emph{Method} (\ancf{coefficients/cert_an.py}, \ancf{coefficients/cert_small_n.py}, \ancf{coefficients/coefflib.py},
\ancf{coefficients/kloost_fast.py}). For each prime power $q\le10^4$ the whole table $j\mapsto S(1,j;q)$ is obtained from one discrete Fourier
transform in Arb, and composite moduli by twisted multiplicativity, $S(1,n;qr)=S(1,n\bar r^2;q)S(1,n\bar q^2;r)$ for $(q,r)=1$ (an elementary consequence of the Chinese
remainder theorem). For prime powers $10^4<q\le6\cdot10^5$ the sums are computed with exact integer residues and a cosine table with a proved
error at most $10^{-14}$, and the rounding of the summation is bounded by Higham's bound \cite{Hig02}; the proved error of each such sum, times its weight in
$\Ss_n$, is below $10^{-16}$. The order derivatives $\Jdot,\Idot$ are evaluated from the series of Lemma~\ref{lem:dot-series} with a geometric
tail bound for arguments $\le8$, and from the closed forms of Lemma~\ref{lem:bessel-basic}(c), with the validated $K_0,K_1$, beyond. The terms
$c>X$ are bounded by the Weil--Estermann bound $|S(\pm1,n;c)|\le d(c)c^{1/2}$ \cite{Wei48,Est61} (see also \cite[Chapter~11]{IK04}), the small-argument
bound of Lemma~\ref{lem:small-w}, and partial summation with $\sum_{c\le u}d(c)\le u(1+\log u)$. The stored balls are rounded outward. The
three coefficient files read by the positivity certificates are in \ancf{coefficients/data/}; a fresh run of the scripts reproduces their
certified numbers (\ancf{coefficients/compare_outputs.py}).

\emph{The trivial-bound box.} The centres $c_n=nT_1(z_n)+\frac14\delta_{n1}$ of \eqref{eq:box} are the terms $c=1$, which \ancf{coefficients/cert_an.py}
computes in Arb at $256$ bits (key \texttt{T1} of \ancf{coefficients/data/an_cert_X1e4.json}); the half-widths $w_n=\abar_n\varrho(z_n)$ are evaluated in
Arb by \ancf{positivity/boxlib.py}. Neither involves a Kloosterman sum with $c\ge2$. In the proofs of the main results the enclosures of this
subsection are used only for the value of $C$ (Appendix~\ref{app:cert-C}); the positivity certificates of the proofs use the box.

\subsection{The window}\label{app:cert-W}

\emph{Statement.} Theorem~\ref{thm:window}, an independent check; script \ancf{positivity/run_window.py} (with \ancf{positivity/cert_window.py} and
\ancf{positivity/poslib.py}).

\emph{Conical series.} For $Y\ge\frac12$ put $z_1:=\sin^2\frac{\theta(Y)}2=1/(4Y^2+1)\le\frac12$, $s:=t^2$ and
$c_k(s):=\prod_{j<k}((\frac12+j)^2+s)/k!^2$. Then $\sqrt{\sin\theta}=2\sqrt{Yz_1}$,
\begin{gather*}
\pc_t(\cos\theta)=\sum_{k\ge0}c_k(s)z_1^k,\\
\pc_t(-\cos\theta)=\frac{\cosh\pi t}\pi\sum_{k\ge0}c_k(s)\big[2\psi(k+1)-2\Re\psi(k+\tfrac12+it)-\log z_1\big]z_1^k
\end{gather*}
by \cite[14.3.1, 15.8.10]{DLMF}, and $1/\pc_t(0)=\Gamma(\frac34+\frac{it}2)\Gamma(\frac34-\frac{it}2)/\sqrt\pi$. All terms of the first series are
positive and increase with $s$, and $c_{k+1}(s)/c_k(s)\le1+s/(k+1)^2$. Both series converge geometrically: the ratio of consecutive terms tends to $z_1\le\frac12$. They are
truncated at $k_{\max}=\lfloor300+8T\rfloor$, $T:=|t_c|+1$, and the tails are bounded with the following lemma, valid uniformly for complex
$t$ with $|t|\le T$.

\begin{lemma}\label{lem:bracket}
Let $T\ge0$, $k_{\max}\ge299+8T$, $k\ge k_{\max}+1$, and $|z|\le T$. Then
$b_k(z):=2\psi(k+1)-\psi(k+\frac12+iz)-\psi(k+\frac12-iz)$ satisfies $|b_k(z)|\le4(\frac12+T)/(k_{\max}+\frac32-T)\le\frac47$.
\end{lemma}

\begin{proof}
On the segments from $k+\frac12\pm iz$ to $k+1$ the real part is at least $k+\frac12-T\ge1$, and $|\psi'(w)|\le\psi'(\Re w)\le2/\Re w$ for
$\Re w\ge1$. Hence $|\psi(k+1)-\psi(k+\frac12\pm iz)|\le2(\frac12+T)/(k+\frac12-T)$; finally $7(\frac12+T)\le k_{\max}+\frac32-T$.
\end{proof}

\emph{Quadrature.} On each panel $[a,b]$ the integrals in \eqref{eq:physical} are replaced by the Gauss--Legendre rule with $m=30$ nodes, whose
nodes and weights are computed as Arb balls; there are $30$ panels on $[\frac12,13]$ for the arcs and $20$ on $[1,13]$ for the axis. The error is
bounded with the following form of \cite[Theorem~19.3]{Tre13}.

\begin{lemma}\label{lem:gauss}
Let $f$ be analytic in the open Bernstein ellipse $E_\rho$ of $[-1,1]$, $\rho>1$, with $|f|\le M$ there, and let $m\ge2$. The $m$-point
Gauss--Legendre rule $G_m$ satisfies $|\int_{-1}^1f-G_m(f)|\le\frac{64}{15}M\rho^{-2m}/(1-\rho^{-2})$.
\end{lemma}

\begin{proof}
The Chebyshev coefficients satisfy $|a_k|\le2M\rho^{-k}$. The rule is exact for polynomials of degree $\le2m-1$, and both the rule and the integral
vanish on the odd $T_k$. For even $k\ge4$, $|\int T_k|=2/(k^2-1)$ and $|G_m(T_k)|\le\sum w_j=2$, so the error on $T_k$ is at most $\frac{32}{15}$; sum
over even $k\ge2m$. (This is \cite[Theorem~19.3]{Tre13} with $n+1=m$ nodes.)
\end{proof}

On each panel, $M$ is bounded on an outward Arb box containing the image of $E_\rho$, uniformly for $t$ in the cell, using $|K_0(z)|\le K_0(\Re z)$,
the absolute conical series with $|z_1|\le1/(4(x_{\rm lo}^2-y_{\rm hi}^2)+1)$, $|\cos(t\log Y)|\le\cosh(t\arctan(y_{\rm hi}/x_{\rm lo}))$, and the
coefficient tail below. The analyticity of the integrand on the box ($\Re Y>0$, $\Re(4Y^2+1)>1.3$, $|\arg Y|+|\arg(4Y^2+1)|<\pi/2$) is checked in Arb.
Beyond $Y=13$: $\int_{13}^\infty\Ax\le\sup_{[13,14]}\Ax/(1-\sqrt{14/13}\,e^{-2\pi})$, which follows from $\Ax(Y+1)\le\sqrt{(Y+1)/Y}\,e^{-2\pi}\Ax(Y)$ (as
$(\log K_0)'\le-1$ and $a_n>0$); $|\Bx|\le\Ax$; and $\ck_t$ is bounded from its series at $z_1=1/677$. The $K_0$-sums use the coefficient balls for
$n\le300$ and $|a_n|\le\abar_n$ with $K_0(x)\le\sqrt{\pi/2x}\,e^{-x}$ and a geometric series beyond.

\emph{Continuum coverage.} On the cell $[t_c-\frac14,t_c+\frac14]$, the quadrature sum $Q$, computed with the truncated conical series, is an analytic
function of $t$ on $|\Im t|<\frac32$; the truncation does not cancel the poles of $1/\pc_t(0)$ at $t=\pm i(\frac32+2m)$. Its Taylor coefficients
at $t_c$ up to degree $19$ are computed by power-series arithmetic in Arb. The Taylor remainder is at most $M_Q(d/r_t)^{20}/(1-d/r_t)$ with
$d=\frac14$, $r_t=1$, and $M_Q\ge\sup_{|z-t_c|=1}|Q(z)|$; $M_Q$ uses $|\pc_z(\cos\alpha)|\le\pc_T(\cos\alpha)$ (Theorem~\ref{thm:kernel-pos}),
$|\cos(z\log Y)|\le\cosh(\log Y)$, and an Arb bound of $1/|\pc_z(0)|$ on $96$ boxes covering the circle. The Taylor polynomial is evaluated on $8$
sub-balls covering the cell, and all error terms are subtracted. Everything is computed at $192$ bits.

\emph{Results} (\ancf{positivity/logs/window_0_40.log}). All $80$ cells are certified; the lower bounds on $[0,\frac12]$ and $[39\frac12,40]$
are $1.12926\cdot10^{-3}$ and $9.034\cdot10^{16}$. Relative to the printed quadrature value $\Pi(t_c)$ at the cell centre, the certified
lower bound is at least $0.2952\,\Pi(t_c)$ on every cell (the smallest ratio, $0.29528$, is on $[\frac12,1]$), the quadrature error is at most
$1.35\cdot10^{-7}\,\Pi(t_c)$ (largest ratio $1.3420\cdot10^{-7}$, on $[39\frac12,40]$) and the Taylor remainder at most $1.33\cdot10^{-7}\,\Pi(t_c)$;
these ratios are informational. The run takes $30$ seconds on $8$ cores. The
enclosure $\Pi(0)\in[1.400011\cdot10^{-3},1.401747\cdot10^{-3}]$ is computed in the same way on a tiny cell around $0$
(\ancf{positivity/check_r0.py}).

\emph{Negative control} (\ancf{positivity/negctl.py}; not used in any proof). With $a_n$ replaced by its term $c=1$, the same certifier proves the
upper bounds of Remark~\ref{rem:negctl} for the right side of \eqref{eq:physical}.

\emph{Independent cross-check} (\ancf{coefficients/cert_Ht.py}, \ancf{coefficients/compare_window.py}; not used in any proof). The
density route of Appendix~\ref{app:cert-C}, run at $s=\frac12+it$, gives certified point values of $H_{\rm raw}(t)$ that do not use the
conical functions or the four-path formula. At all $16$ cell centres in $[0,8]$ and at $t=10.25,15.25,20.25,30.25,39.75$, these enclosures of
$\Pi(t)=2\pi^2(t^2+\frac14)^2H_{\rm raw}(t)$ overlap those of the window certificate ($21$ of $21$). The two routes share only the coefficient
enclosures.

\subsection{The moment certificate}\label{app:cert-M}

\emph{Statement.} Theorem~\ref{thm:moments}; script \ancf{positivity/cert_moments_box.py} (with \ancf{positivity/cert_moments.py} and
\ancf{positivity/boxlib.py}). The same computation with the certified coefficients, \ancf{positivity/cert_moments.py}, gives the check $|t|\le10.402$
and the numbers of Remarks~\ref{rem:sign} and~\ref{rem:moments}.

\emph{The transform at a node.} For real $x>1$, $\Gh_{\rm raw}(x)=\sin^2(\pi x)\sqrt x\,[\sum_{n\le N}a_nK_0(4\pi\sqrt{nx})+T_N(x)]$. The tail $T_N$ is
bracketed in closed form from the bounds $\abar_n(\lambda_0(z_n)-\eta_N)\le a_n\le\abar_n(1-1/z_n+\eps_N)$ for $n>N$, which hold on $\TB$; here
$\eta_N$ and $\eps_N$ are the remainder terms of $\lambda_-$ and $\lambda_+$ at $z_{N+1}$, and $\lambda_0(z)\ge1-\frac{11}{8z}-(\frac3{32}+\frac{\pi^3}{16})z^{-2}$. With
$\sqrt{\pi/2y}\,e^{-y}(1-\frac1{8y})\le K_0(y)\le\sqrt{\pi/2y}\,e^{-y}$ the sums over $n>N$ become sums of $e^{-\kappa\sqrt n}$ and
$n^{-1/2}e^{-\kappa\sqrt n}$, $\kappa=4\pi(\sqrt x-1)$, which are compared with integrals. $N$ is chosen per node: $60$ near $x=1$ for $\TB$
($1500$ for the certified coefficients), and about $20$ for $x\ge10$.

\emph{Quadrature.} The integrals \eqref{eq:moments} are split at integers, where $W_j$ is not analytic, and computed by $24$-node
Gauss--Legendre rules on $99$ panels: $[1+2^{-k-1},1+2^{-k}]$ for $1\le k\le21$ ($\rho=2$), $[1.5,1.75]$ and $[1.75,2]$ ($\rho=3$), and two panels in each
$[k,k+1]$, $2\le k\le39$ ($\rho=4$). The error is bounded by Lemma~\ref{lem:gauss}, with $M$ an Arb bound on an outward box containing the
Bernstein ellipse, from $|K_0(w)|\le K_0(\Re w)$, $\Re\sqrt w\ge\sqrt{\Re w}$, $|\sin\pi w|\le\cosh(\pi\Im w)$ and $|\sin\pi w|\le\sinh(\pi|w-k|)$; only
real arguments of $K_0$ are evaluated.

\emph{Ends.} Every $(b_n)\in\TB$ satisfies $0<b_n\le\abar_n$, so the transform formed from it lies between $0$ and the function $U$ of the last
paragraph below. On
$[1,1+2^{-22}]$ this gives $\Gh_{\rm raw}(x)\le x^{1/4}[(x-1)^2/16+x/(32\pi^2)]$, by $\sin^2(\pi x)\le\pi^2(x-1)^2$, $\kappa\ge2\pi(x-1)/\sqrt x$ and
$e^{-\kappa}(1+\kappa)\le1$. On $[40,\infty)$ it gives $\Gh_{\rm raw}(x)\le x^{1/4}e^{-\kappa}(1+2(1+\kappa)/\kappa^2)/(16\pi^2)$; with
$\sum_{m\le x}d(m)m^{-1/2}\le x(1+\log x)$ the integral is bounded by an incomplete Gamma function. For $M_0$ the quadrature error, the first end and the second end are at most $2.1\cdot10^{-17}$, $2.4\cdot10^{-10}$ and
$6\cdot10^{-31}$.

\emph{Uniformity on the box.} $P_7(t)$ is linear in the coefficient sequence. Let $\pi_n(t)$ be its value for the sequence that is $1$ at $n$ and
$0$ elsewhere. For every $b\in\TB$ the certificate uses
\[
P_7(t;b)\ge P_7^{\rm low}(t)-\sum_{n\le8}w_n|\pi_n(t)|-\eps\,P_7^{\rm abs}(t),
\]
where $P_7^{\rm low}(t)$ is the lower bound obtained from the centres $c_n$ for $n\le60$ and the brackets above for $n>60$,
$\eps:=\max_{8<n\le60}w_n/(c_n-w_n)=1.7331\cdot10^{-8}$, and $P_7^{\rm abs}(t)$ is the same polynomial with all signs $+$ and the upper ends of the
moments of the centres.

\emph{Continuum in $t$.} On cells $[t_a,t_b]$ of width $1/256$, the even and odd parts of $P_7$ are increasing in $t\ge0$, because the moments are
non-negative; so $P_7\ge P_+(t_a)-P_-(t_b)$, where $P_+$ uses the lower ends of the even-indexed terms and $P_-$ the upper ends of the
odd-indexed ones. No Taylor model in $t$ is needed.

\emph{Results} (\ancf{positivity/logs/cert_moments_box.log}; $64$ bits, $14$ seconds on $6$ cores). The lower bound is positive on every cell of
$[0,10.35546875]$, and the verdict line checks that this interval contains $[0,9]$. With the certified coefficients
(\ancf{positivity/logs/cert_moments_true.log}; $128$ bits, $20$ seconds), $M_0=0.00028062\pm4.3\cdot10^{-9}$ and the bound is positive on
$[0,10.40234375]$; this log also prints the enclosure of $M_0/\Xi(0)^2$ quoted in Remark~\ref{rem:sign} and the ratios of Remark~\ref{rem:moments}.

\emph{The elementary range} (\ancf{positivity/cert_elementary.py}; not used in any proof). From Theorem~\ref{thm:an-bounds} and the bounds for $K_0$
above, $L\le\Gh_{\rm raw}\le U$ on $(1,\infty)$, where
\[
U(x)=\frac{\sin^2(\pi x)\,x^{1/4}e^{-\kappa}}{16\pi^2}\Big(1+\frac{2(1+\kappa)}{\kappa^2}\Big),\qquad
L(x)=\frac{\sin^2(\pi x)\,x^{1/4}}{16\pi^2}\Big(\frac{2e^{-\kappa}(1+\kappa)}{\kappa^2}-\frac{2.125}{2\pi\kappa}\Big).
\]
Arb's rigorous integration gives $M_0\ge\frac1\pi\int_1^{1.3}L\,\frac{\dd x}x\ge1.6932\cdot10^{-4}$ and
$M_2\le\frac1\pi\int_1^\infty W_1U\,\frac{\dd x}x\le4.2931\cdot10^{-7}$, so $P_0(t)=M_0-2\pi^2t^2M_2>0$ for $|t|<4.4699$. Keeping the term $n=1$, with
$a_1\in[c_1-w_1,c_1+w_1]$, gives $M_0\ge1.9672\cdot10^{-4}$, $M_2\le3.6864\cdot10^{-7}$ and $|t|<5.1995$ (\ancf{positivity/logs/cert_elementary.log}).

\subsection{The large-\texorpdfstring{$|t|$}{|t|} numbers}\label{app:cert-L}

\emph{Statement.} Propositions~\ref{prop:L-numbers-box} and~\ref{prop:L-numbers}; scripts \ancf{positivity/cert_large_t.py} (items (i), (iii), (iv))
and \ancf{positivity/cert_gtail.py} (item (ii)). For Proposition~\ref{prop:L-numbers-box} the wrapper \ancf{positivity/run_on_box.py} runs the same two
scripts, unchanged, with the coefficient balls for $n\le300$ replaced by the intervals of $\TB$ (\ancf{positivity/boxlib.py}); beyond $n=300$ the
scripts use $0<a_n\le\abar_n$, which holds on $\TB$. The scripts enclose every quantity they bound in Arb as a function of the coefficient balls,
so the certified bounds hold for every sequence in $\TB$.

\emph{Method.} With $(T_0,Y_z)=(9,0.90)$ or $(8,0.870)$: $p_+(T_0)$ is bounded below by a Riemann sum over $2000$ geometric cells of $[1,4]$, with the
integrand enclosed on each cell and the sum rounded down; $p_-(T_0)$ is bounded above by a Riemann sum over $8000$ cells of $[\frac12,Y_z]$, with
$\Bx^-$ bounded by the negative of the lower end of the ball of $\Bx$, rounded up; $\pc_0(\pm\cos\theta)$ come from the conical series at $t=0$. The
cell minima of $\Bx$ in (i) are the minima of the exact lower ends over $300$ cells of $[Y_z,1]$ and the $2000$ cells of $[1,4]$. $I_{\rm ax}$ is computed by Gauss--Legendre
quadrature with Lemma~\ref{lem:gauss} (quadrature error about $2.5\cdot10^{-17}$, tail beyond $Y=13$ about $5\cdot10^{-32}$, both included in the bound). For (ii), the terms
$2\le n\le399$ of $\bar\rho(4)$ are summed in Arb with $a_1$ replaced by the lower end of its ball, which gives an upper bound for $\bar\rho(4)$, and the terms $n\ge400$ are at most twice the
term $n=400$, since consecutive ratios are at most $(1+\frac1n)e^{2\pi/\sqrt n-8\pi}<10^{-10}$. The lower bound in (iv) is formed in Arb from the exact
endpoints of the bounds in (iii). The runs take $133$ seconds (Proposition~\ref{prop:L-numbers-box}, \ancf{positivity/logs/cert_large_t_T9_box.log}) and
$138$ seconds (Proposition~\ref{prop:L-numbers}) on one core.

\subsection{The normalising constant}\label{app:cert-C}

\emph{Statement.} Proposition~\ref{prop:C}; script \ancf{coefficients/cert_H0.py}.

\emph{The formula.} Put $g_{\rm raw}:=\Gh_{\rm raw}/(2\pi\sqrt x)$. By \eqref{eq:mellin-H}, $H_{\rm raw}(0)=\int_0^\infty g_{\rm raw}(x)x^{-1/2}\dd x/\gi(0)^2$ with
$\gi(0)^2=\pi^{-1/2}\Gamma(\frac14)^2/64$. Let $d:=-D/(16\pi^2)=v^{-5/2}\sum_n2na_nK_0(2\pi n/v)$, $d_{\rm grow}:=-D_{\rm grow}/(16\pi^2)$ and
$d_{\rm rest}:=d-d_{\rm grow}$. By Theorem~\ref{thm:lift}(d) and the integrals in its proof, $\Gh_{\rm raw}(x)=\sin^2(\pi x)\sqrt x\,[P_L(x)+\int_0^\infty
d_{\rm rest}(v)\sqrt vK_0(2\pi vx)\dd v]$ with $P_L:=-R_L/(16\pi^2)$, that is, $P_L(x)=\frac1{32\pi^4(x^2-1)}+\frac1{8\pi^4(x^2-1)^2}$. Fubini's theorem gives
\begin{gather*}
2\pi\int_0^\infty g_{\rm raw}(x)x^{-1/2}\dd x=\int_0^\infty\sin^2(\pi x)x^{-1/2}P_L(x)\dd x+\int_0^\infty d_{\rm rest}(v)\sqrt v\,\mathcal Q(v)\dd v,\\
\mathcal Q(v)=\frac{\Gamma(\frac14)^2\,[1-{}_2F_1(\frac14,\frac14;\frac12;-v^{-2})]}{2^{5/2}(2\pi v)^{1/2}},
\end{gather*}
where $\mathcal Q(v)=\int_0^\infty\sin^2(\pi x)x^{-1/2}K_0(2\pi vx)\dd x$ is evaluated by \cite[6.699.12]{GR15}. On $(0,1]$, $d_{\rm rest}$ is computed from
the series for $d$ and the closed form of $d_{\rm grow}$; on $[1,8]$ from the following expansion at the cusp $\infty$, with the coefficients
$A_n(1)$ and $A_n'(1)=4\Ss_n$ of Theorem~\ref{thm:niebur-fourier}.

\begin{lemma}\label{lem:D-infinity}
For $v>0$, with $y_n:=2\pi nv$,
\begin{multline*}
D(v)=D_{\rm grow}(v)+4\pi v^{3/2}\big[K_1(y_1)-y_1K_0(y_1)\big]\\
+\sum_{n\ge1}4\pi nv^{3/2}\Big[A_n'(1)\big(K_1(y_n)-y_nK_0(y_n)\big)+A_n(1)\Big(\frac{K_0(y_n)}{y_n}-2K_1(y_n)\Big)\Big].
\end{multline*}
\end{lemma}

\begin{proof}
Write $\dot P=\sum_np_n(v)\sin2\pi nu$ and $P_1=\sum_nr_n(v)\sin2\pi nu$ as in Proposition~\ref{prop:sym2}(c). From \eqref{eq:Phi},
$D=-E_1(iv)-\partial_uE_0(iv)=\sum_n2\pi n[3vp_n+2v^2p_n'-2vr_n]$. Insert $p_n,r_n$ and use $K_0'=-K_1$, $(yK_1)'=-yK_0$, $I_0'=I_1$; the terms
$-\frac{I_0(2\pi v)}{2\pi\sqrt v}$ of $p_1$ and $\sqrt vI_1(2\pi v)$ of $r_1$ give $D_{\rm grow}$.
\end{proof}

The two expressions for $d_{\rm rest}$ give overlapping enclosures at $v=1$. The integrals are computed with Arb's rigorous integration routine,
with $K_\nu$ of complex argument evaluated at exact midpoints plus a derivative bound for the radius; the tails, $n>60$ at the cusp $0$
and $n>30$ at the cusp $\infty$, use
$|a_n|\le\abar_n$, $|A_n'(1)|\le4\abar_n/n$ and Lemma~\ref{lem:alpha}. The ends are bounded analytically: on $[0,10^{-8}]$, $0\le d\le10^{-6}\sqrt v/(8\pi^2)$
and $|\mathcal Q(v)|\le2^{-3/2}\Gamma(\frac14)^2(2\pi v)^{-1/2}$ give a contribution at most $3.2\cdot10^{-14}$; on $[8,\infty)$ a termwise majorant from
Lemma~\ref{lem:D-infinity}, $K_0\le K_1\le K_{3/2}$ and $|\mathcal Q(v)|\le\sqrt2\pi^2\Gamma(\frac54)^2(2\pi v)^{-5/2}$ gives at most $3.5\cdot10^{-20}$; and
the part $x>2000$ of the first integral is bracketed with the second mean value theorem.

\emph{Results} (\ancf{coefficients/logs/cert_H0.log}, $160$ bits). With the enclosures of $a_1,a_2$ for $c\le6\cdot10^5$: the first integral is
$0.0034380158943\pm2.9\cdot10^{-13}$; the second is $-0.0152642540638\pm7.7\cdot10^{-10}$ on $(0,1]$ and $0.0126529953780\pm5.8\cdot10^{-10}$ on $[1,\infty)$;
so $2\pi\int g_{\rm raw}x^{-1/2}=0.00082675720847\pm1.4\cdot10^{-9}$ and $H_{\rm raw}(0)=0.0011355097904\pm1.86\cdot10^{-9}$. The radius is dominated by the
bound for the terms $c>6\cdot10^5$ of $\Ss_1$. A second run that splits at $v=1.25$ instead of $v=1$ gives an overlapping enclosure. The run takes
about $90$ seconds.

\subsection{The constants of Section~\ref{sec:fourier}}

The constants $\lambda_0(4\pi)$, $\varrho(4\pi)$, $\lambda_\pm(4\pi)$, the two monotone envelopes in the proof of Theorem~\ref{thm:an-bounds} and
$\varrho_\alpha(4\pi)$ are elementary expressions evaluated at one point; \ancf{coefficients/analytic_constants.py} evaluates them in Arb, for
example $\lambda_-(4\pi)=0.850104915628\pm2\cdot10^{-13}$.

\subsection{The Lean formalisation}\label{app:lean}

The ancillary directory \ancf{lean/} contains a Lean~4 project, built on Mathlib and on the corresponding project for Part~I, that formalises
the logical spine of the proof. It states Theorems~\ref{thm:main-S} and~\ref{thm:main-U}, Corollaries~\ref{cor:main-RH} and~\ref{cor:nogap} and the reduction
(Lemma~\ref{lem:voronoi}(a) and Proposition~\ref{prop:reduction}), and derives them from Mathlib, from six axioms of Part~I's ledger, and from
exactly two axioms for this paper. The six axioms of Part~I are the explicit formula (Lemma~\ref{lem:I-EF}) and the bound \eqref{eq:XB}, used for
Theorem~\ref{thm:main-S}; the zero-side support theorem (Theorem~\ref{thm:I-zerosupport}) and Theorem~\ref{thm:I-logic}(b), added for
Theorem~\ref{thm:main-U}; the duality theorem (Theorem~\ref{thm:I-duality}), added for Corollaries~\ref{cor:main-RH} and~\ref{cor:nogap}; and the lower bound for
$\kstar$ of Part~I, Proposition~2.8, used only to show that $\kstar$ and $\kOPS$ are finite. The Part~I axiom for
Theorem~\ref{thm:I-zerosupport} states only its first conclusion, $\nu=\nuz$; the remaining conclusions, RH and $\mu=\muz$, are derived in Lean
from Theorem~\ref{thm:I-logic}(b), as in Part~I. The two axioms for this paper are: first, one axiom stating
parts (b) and (c) of the Voronoi decoupling (Lemma~\ref{lem:voronoi}); second, one joint existence axiom for the object: an $H\in\Wc_\delta$, for
some $\delta<\frac12$, with (C3), $\Gh_H\ge0$ on $[0,\infty)$, $\Gh_H>0$ on $[0,\xin2)$ and $H>0$ on $\RR$, which is what
Theorem~\ref{thm:main-object}, Corollary~\ref{cor:Pi-positive} and the proof of Corollary~\ref{cor:nogap}(a) provide. The
project builds with no unproved step (\texttt{sorry}) and passes its audit script; its \ancf{README.md}, \ancf{LEDGER.md} and \ancf{FAITHFUL.md}
describe the axioms and list every difference between the Lean statements and the paper's. In particular the Lean form of the
first part of Theorem~\ref{thm:main-S} asserts less about $H$ than the paper does: membership in $\Wc_\delta$ for one $\delta<\frac12$, that is, the
exponent $8$ of (C1) instead of $9-\delta$ for every $\delta<1$, and no entireness. Moreover the differentiability of $\hF$ is not formalised:
the Lean form of $\hF'(\xin n)=0$ says that the derivative of $\hF$ at $\xin n$, if it exists (\texttt{HasDerivAt}), is $0$. The Lean kernel
therefore checks the assembly of the proof, namely that the main statements follow from these inputs. It
does not check the inputs: Sections~\ref{sec:automorphic}--\ref{sec:physical}, the appendices and the certificates are not formalised.
Corollary~\ref{cor:nogap} is formalised as \texttt{rh\_iff\_kappaOPS\_zero} (with \texttt{corollary8\_2}, \texttt{corollary8\_2\_b},
\texttt{FT\_window}; Lemma~\ref{lem:transfer} as \texttt{transfer}), except the zero set in (a) and clause (c)(iv).
